% ERRORQUAKE Theoretical Framework — New Section 3
% Non-Reducibility Theorem, Resolution Bound, Multiplicative Error Model
%
% To be integrated into main.tex as Section 3.

\subsection{Theoretical Framework}
\label{sec:theory}

We establish three formal results that underpin the empirical analysis.
Theorem~\ref{thm:nonreduce} proves that severity distributions carry
information not reducible to error rate; Proposition~\ref{prop:resolution}
bounds the measurement precision needed to detect this information;
and Proposition~\ref{prop:multiplicative} connects the observed distribution
families to a mechanistic model of error generation.

\begin{theorem}[Non-Reducibility of Severity Profile]
\label{thm:nonreduce}
Let $\mathcal{M}$ be a set of language models, each characterised by an
error rate $\varepsilon_M = P(S_M > 0)$ and a severity distribution
$F_M(s) = P(S_M \leq s \mid S_M > 0)$ on a bounded support $[0, s_{\max}]$.
Define the severity distribution index $b_M$ as the slope of the
Gutenberg--Richter fit $\log_{10} N(M \geq m) = a - b(m - m_{\min})$.
Then:

\begin{enumerate}[label=(\roman*)]
\item \textbf{Existence of matched-accuracy divergence.} For any
$\delta > 0$, there exist models $M_i, M_j$ such that
$|\varepsilon_i - \varepsilon_j| < \delta$ while
$|b_i - b_j|$ is arbitrarily large (up to the support-imposed maximum).

\item \textbf{Informational non-redundancy.} The mutual information
$I(b; \text{model} \mid \varepsilon) > 0$ whenever the population of
models includes pairs with matched accuracy and divergent severity
profiles.
\end{enumerate}
\end{theorem}

\begin{proof}[Proof sketch]
Part (i) is constructive. Consider two Gutenberg--Richter distributions
on the discrete grid $\{0.5, 1.0, \ldots, 4.0\}$ with slopes $b_1$ and
$b_2 \neq b_1$. For any target error rate $\varepsilon^*$, the error rate
of the $j$-th model is $\varepsilon_j = \sum_{m > 0} P_j(M = m)$, which
is a monotone function of the normalisation constant $a_j$ for fixed $b_j$.
Adjusting $a_j$ to satisfy $\varepsilon_j = \varepsilon^*$ while keeping
$b_j$ fixed produces two models with identical error rate and distinct
severity profiles.

Part (ii) follows from the data processing inequality: if $b$ is a
deterministic function of $F$ and $\varepsilon$ is a different deterministic
function of $F$, then $I(b; \text{model} \mid \varepsilon) \geq 0$ with
equality only when $b$ is a deterministic function of $\varepsilon$. The
constructive example in (i) provides a counterexample to this functional
dependence.

\emph{Empirical confirmation}: on our 21-model catalog, $I(b; \text{model}
\mid \varepsilon) = 1.56$ bits (5-bin discretisation), and only 35.5\%
of cross-model $b$ variance is explained by $\varepsilon$ ($R^2 = 0.356$).
The remaining 64.5\% is independent severity information.
\end{proof}

\begin{proposition}[Resolution Bound]
\label{prop:resolution}
Let $b$ be estimated from $n$ severity scores on a $k$-level ordinal
scale with inter-rater reliability $\text{ICC}(2,k) = r$. The standard
error of $\hat{b}$ satisfies
\begin{equation}
\text{SE}(\hat{b}) \geq \frac{c(b, k)}{\sqrt{n \cdot r}}
\label{eq:resolution}
\end{equation}
where $c(b, k)$ is a constant depending on the true $b$-value and the
number of scale levels. Two models with $b$-values $b_1, b_2$ are
distinguishable at significance $\alpha$ with power $1 - \beta$ when
\begin{equation}
|b_1 - b_2| \geq (z_{\alpha/2} + z_\beta) \cdot \sqrt{2} \cdot \text{SE}(\hat{b}).
\label{eq:min_delta}
\end{equation}
\end{proposition}

\begin{proof}[Proof sketch]
The maximum-likelihood estimator $\hat{b}$ from the Aki formula has
asymptotic variance $\text{Var}(\hat{b}) = b^2 / n_{\geq m_{\min}}$
\citep{aki1965}. When scores are measured with reliability $r < 1$,
the effective sample size is attenuated: measurement noise inflates
the variance by a factor $1/r$, giving
$\text{Var}(\hat{b}) \geq b^2 / (n_{\geq m_{\min}} \cdot r)$.
The two-sample comparison follows from the standard normal approximation.

\emph{Application to ERRORQUAKE}: with median $\text{SE}(\hat{b}) = 0.064$,
the minimum detectable $\Delta b$ at $\alpha = 0.05$, power $= 0.80$ is
$0.253$. The observed $b$ range across models is $[0.57, 1.31]$
(spread $= 0.74$), confirming the study is powered to detect the reported
differences despite the moderate ICC.
\end{proof}

\begin{proposition}[Multiplicative Error Model]
\label{prop:multiplicative}
If the severity of an error is generated by a multiplicative process
$S = \prod_{l=1}^{L} X_l$ where $X_l$ are independent positive random
variables representing the contribution of processing stage $l$
(e.g., retrieval, reasoning, generation), then by the central limit
theorem for products, $\log S \xrightarrow{d} \mathcal{N}(\mu_L, \sigma_L^2)$
as $L \to \infty$, i.e., $S$ is approximately lognormal.

Furthermore, if $\sigma_l^2 = \sigma_0^2$ for all $l$, then
$\sigma_L^2 = L \sigma_0^2$ and the Gutenberg--Richter slope satisfies
$b \propto 1 / \sigma_L \propto 1 / \sqrt{L}$: models with more
multiplicative stages (more layers) have smaller $b$ (heavier tails).
\end{proposition}

\begin{proof}[Proof sketch]
Taking logarithms, $\log S = \sum_{l=1}^{L} \log X_l$. If the
$\log X_l$ have finite mean $\mu_0$ and variance $\sigma_0^2$, the
classical CLT gives $\log S \sim \mathcal{N}(L\mu_0, L\sigma_0^2)$
for large $L$. The Gutenberg--Richter slope of a lognormal is
$b = \log_{10}(e) / \sigma_L$, decreasing in $\sigma_L$.

\emph{Empirical support}: 10/21 models in our catalog have stretched
exponential (a generalisation of lognormal) as the BIC-best fit, and
2/21 have lognormal. No model is best-fit by a pure power law, consistent
with the finite-$L$ multiplicative prediction rather than an
infinite-variance process. The negative scaling correlation
($\rho_s = -0.56$, $n_{\text{dense}} = 14$) is directionally consistent
with the $b \propto 1/\sqrt{L}$ prediction, though the study is not
powered to test the functional form.
\end{proof}
